\documentclass[11pt]{article}
\usepackage[margin=1in]{geometry}
\usepackage[T1]{fontenc}
\usepackage[utf8]{inputenc}
\usepackage{amsmath,amssymb,amsthm,mathrsfs}
\usepackage[hidelinks]{hyperref}
\newtheorem{theorem}{Theorem}
\newtheorem{lemma}[theorem]{Lemma}
\newtheorem{corollary}[theorem]{Corollary}
\newcommand{\DD}{d'd''}
\newcommand{\PH}{\operatorname{PH}}
\newcommand{\MA}{\operatorname{MA}}
\title{Bounded continuous Monge--Ampere uniqueness\\by coordinate curves and Gauss rays}
\author{Working paper; independent review pending}
\date{9 October 2026}
\begin{document}
\maketitle
\begin{abstract}
For the specified compact analytic torus quotient, two bounded
continuous potentials in the entire minimal continuous CLD psh sheaf
with the same actual Monge--Ampere measure differ by a constant.
We include the bounded Green argument and divisor testing, then prove
a bounded harmonic Liouville lemma for the analytic multiplicative
group over finite extensions of the actual base field. Finite Galois
descent of signed Radon Laplacians, calibrated Gauss-ray integration,
and density of closed rigid points suffice. Coordinate algebraic
curves then connect the values of the difference; no recognition by
all analytic curves after arbitrary scalar extensions is assumed.
All claims are working claims pending independent review. The result
does not settle unbounded finite-energy uniqueness, finite-energy
compactness, obstacle orthogonality, or general equation existence.
\end{abstract}

\section{Carrier, domain, and exact foundations}
Put $K=\mathbb C((s))$, $|s|=e^{-1/8}$, and let
\[
 X=((\mathbb G_m)^2)^{\mathrm{an}}/
 \langle(s^8,2),(3,s^8)\rangle,
 \qquad x_i=-\log|z_i|,
 \quad \rho=\tfrac12|x|^2,
 \quad\theta=\DD\rho.
\]
We retain the compact, connected, smooth, good, boundaryless analytic
quotient and its ordinary injective torus-chart atlas as geometric
foundations. Tropical translation by the periods is $\mathbb Z^2$;
rectangles of coordinate width less than one give injective charts.
The local reference potentials differ by affine functions of unit
logarithms, so their Hessians glue. Quotient existence, separatedness,
properness and nonalgebraization are separate geometric statements;
the present theorem does not reestablish their entire proof trees.

Let $\mathscr P$ be the minimal sheaf of finite continuous functions
containing ordinary smooth CLD psh functions and closed under locally
uniform limits. Let
\[
 \mathcal C=\{u\in C(X,\mathbb R):\rho+u\in\mathscr P
                         \text{ on every chart}\}.
\]
Here every continuous-sheaf stage is retained. Convex combinations and
finite maxima preserve this class: prove the assertions for smooth
weights by convexity and smooth convex soft maxima, then use minimality
and uniform closure successively in the two slots.

The immediately applied CLD v2 foundations are Definition8.1.7 and
Lemma8.1.8; Lemmas8.5.1,8.5.3,8.5.5; the symmetric recursive continuous
products in8.7.1,8.7.2,8.7.4; multilinear continuous dpsh calculus and
smooth null changes in8.8.4--5; scalar Radon recognition and
smooth-coefficient measures in7.2.8,7.2.11--12; compact-support
Green/Stokes in5.11.1; calibrated integration15.7.1; and
Poincare--Lelong16.7.3~\cite{CLD}.
They give positive closed products, local uniform continuity, and the
following multiplier rule: if positive currents $T_j$ converge weakly
to $T$ and continuous coefficients converge locally uniformly, their
products converge weakly. We use only scalar and surface $(1,1)$
positivity; no higher-degree strong-cone order unit is assumed.
Smooth cutoffs and the underlying calibrated integration are ordinary
published foundations. The whole immediate applied source units were
read; deeper recursive proofs are not all newly audited.

We also use the defining Tate-algebra presentation of a rational domain
and the density of finite polynomials in a completed Tate algebra;
the latter follows directly by truncating its convergent series.
The ordinary presentation and its universal-property argument are
Van der Put2.1--2.2~\cite{VdP}; the strict presentation is also stated
in Papikian3.4--3.5~\cite{Pap}.
This includes generalized polyannulus radii. Rational domains form a
compact neighborhood basis inside a good affinoid torus chart.
These are ordinary affinoid foundations, not a claimed approximation
by polynomials on an arbitrary open set without shrinking.

\section{Bounded Green calculus at all continuous stages}
Fix a positive closed $(1,1)$ current $T$. Complete compact smooth
$(1,0)$ forms modulo the null space of the positive pairing
$\int a\wedge Jb\wedge T$. With a nonnegative compact smooth cutoff
$\chi$, write the local gradient pairing as $Q_{T,\chi}$.
Let $\mathscr H_T$ consist of continuous functions locally approximable
uniformly by smooth functions whose gradients converge strongly in
these cutoff seminorms.

Gradients are unique: if smooth $f_j\to0$ uniformly and their gradients
converge strongly to $a$, then, for each compact smooth form $b$,
closedness of $T$ gives
\[
 \int d'f_j\wedge Jb\wedge T
 =-\int f_jd'(Jb)\wedge T\longrightarrow0.
\]
The right side pairs $f_j$ with a finite signed measure. Smooth forms
are dense in the completion, so $a=0$. Smaller patches and cutoffs
give the same local assertion. Smooth partitions patch approximants:
the gradient of $\sum\chi_i f_{i,j}$ splits into
$\sum\chi_i d'f_{i,j}+\sum f_{i,j}d'\chi_i$, with both sums strongly
convergent. Bounded continuous multiplication of forms follows by
scalar smooth approximation and the norm bound. Products and cutoffs
of functions in $\mathscr H_T$ therefore have the usual gradient rule.

If $f\in\mathscr H_T$ and $L_Tf=\DD(fT)$ is signed Radon, smooth
approximation first in $f$, then uniformly and strongly in a compactly
supported $g\in\mathscr H_T$, proves
\[
 \int gL_Tf=-Q_T(f,g).
\]
In particular the product rule yields
\[
 Q_{T,\chi}(w,w)
 =\tfrac12\int w^2\DD\chi\wedge T-\int\chi w L_Tw.
\]
The sheaf of $f\in\mathscr H_T$ with $L_Tf\ge0$ contains smooth psh
functions and is uniformly closed. Indeed choose nested cutoffs
$\chi,\eta$, with $\eta=1$ near $\operatorname{supp}\chi$, and a common
bound $|f_j|\le B$ on the whole support of $\eta$. For
$\epsilon=\|f_j-f_k\|$ there, positivity and integration by parts give
\[
 0\le Q_{T,\chi}(f_j-f_k,f_j-f_k)
 \le\tfrac{\epsilon^2}{2}\|\DD\chi\wedge T\|_{TV}
 +2B\epsilon\|\chi\|_\infty\|\DD\eta\wedge T\|_{TV}.
\]
To obtain smooth approximants to the limit on a smaller patch, choose
an approximant to each $f_j$ with uniform and gradient error at most
$1/j$ for a cutoff equal to one on that patch. The displayed estimate
proves strong convergence. Positivity survives distributional limits.
Minimality now puts all of $\mathscr P$ in this sheaf. Its vector span
contains continuous dpsh functions. Finite patching on compact $X$
proves the global formula
\[
 \int f\DD(gT)=\int g\DD(fT)=-Q_T(f,g),
 \qquad Q_T(f,f)\ge0,
\]
and Cauchy--Schwarz for continuous dpsh functions. This argument uses
positive Laplacians and signed measures, not arbitrary continuous
multiplier closedness.

For $u\in\mathcal C$ write $\theta_u=\theta+\DD u$ and
$B(u,v)=\theta_u\wedge\theta_v$ for the actual recursive local
products, glued through the smooth null potential changes.
Set $\MA(u)=B(u,u)$. Every mixed measure has the same finite mass
$M=\int\theta^2=2$; Green/Stokes annihilates the exact terms.

\section{An elementary bounded all-test estimate}
Let $a,b,w\in\mathcal C$, set $f=a-b$, $g=(a+b)/2$, $h=w-g$, and
\[
 I(a,b)=\int(a-b)(\MA(b)-\MA(a)).
\]
The bounded product expansion gives
\[
 I=Q_{\theta_a}(f,f)+Q_{\theta_b}(f,f)\ge0.
\]
Multilinearity and Green give the signed identity
\[
 Q_{\theta_w}(f,f)-Q_{\theta_g}(f,f)
 =Q_{\theta_a}(f,h)-Q_{\theta_b}(f,h).
\]
Indeed both sides are $-\int f\DD f\wedge\DD h$.
Thus Cauchy--Schwarz implies
\[
 Q_{\theta_w}(f,f)\le I/2+
 \sqrt{I\bigl(Q_{\theta_a}(h,h)+Q_{\theta_b}(h,h)\bigr)}.
\]
For an explicit bounded version choose $B_0$ with
$|a|,|b|,|w|\le B_0$. Then $|h|\le2B_0$, and
\[
 Q_{\theta_a}(h,h)
 =-\int h\,[B(a,w)-(\MA(a)+B(a,b))/2]
 \le4MB_0.
\]
The same holds in the $b$ slot, since every positive measure in the
brackets has mass $M$. Consequently
\[
 Q_{\theta_w}(a-b,a-b)\le I/2+\sqrt{8MB_0I}.
\]
This estimate is sufficient here; it claims no uniform sup-norm bound
for unbounded generated presentations.

\begin{lemma}\label{annihilate}
If $u,v\in\mathcal C$ satisfy $\MA(u)=\MA(v)$ and $d=u-v$, then
$\DD(d\theta_w)=0$ for every $w\in\mathcal C$.
\end{lemma}
\begin{proof}
Equality of the measures gives $I(u,v)=0$. The estimate gives
$Q_{\theta_w}(d,d)=0$. Cauchy--Schwarz with compact smooth scalar
tests and the Green identity show that every test pairing of
$\DD(d\theta_w)$ vanishes. The difference is bounded and continuous.
\end{proof}

\section{Laurent divisor currents from glued smooth caps}
Take a nonzero Laurent polynomial $A$. At a point of its zero locus
choose an injective chart $|x_i-a_i|<r<1/2$, and put
$\rho_a=|x-a|^2/2$. Choose a smooth convex $\chi$ with
$\chi=0$ on $(-\infty,0]$, $0\le\chi'\le1$, and
$\chi(t)=t-c$ for $t\ge1$. The functions
\[
 \psi_N=\chi(\log|A|+N)
\]
extend smoothly by zero near the zeros of $A$ and are ordinary smooth
psh: away from those zeros $A$ is a unit moment, and at each zero the
cap is locally constant.

Fix $0<\delta<r^2/4$ and choose $\epsilon_N>0$ such that
$\delta+\epsilon_N\sup\psi_N<r^2/2$ on the closed polyannulus.
Then
\[
 w_N=\max(\rho_a,\delta+\epsilon_N\psi_N)-\rho_a
\]
extends by zero to a member of $\mathcal C$. The strict boundary
inequality supplies a collar where it is zero, and smooth soft maxima
give the required locally uniform psh approximation. On the fixed
smaller neighborhood $\rho_a<\delta/2$ one has
$\theta_{w_N}=\epsilon_N\DD\psi_N$. Lemma\ref{annihilate} gives
$\DD(d\DD\psi_N)=0$ there.

CLD16.7.3 applies to the regular holomorphic $A$ on the smooth,
pure-dimensional, boundaryless torus chart:
\[
 \DD[\log|A|]=\delta_{\operatorname{div}A}.
\]
The lower-dimensional support fact used in its immediate proof says
that compact smooth top-degree tests avoid $\{A=0\}$.
On such a test support, $\psi_N-N=\log|A|-c$ for all sufficiently large
$N$. Thus the positive currents $\DD\psi_N$ converge weakly to the
actual divisor current. CLD8.5.1 with the continuous coefficient $d$,
followed by continuous differentiation of currents, proves
\[
 \DD(d\delta_{\operatorname{div}A})=0
\]
on the smaller neighborhood. Covering the divisor proves the local
identity everywhere in every chart; off the divisor it is automatic.
The argument does not assert closedness of $d\delta_{\operatorname{div}A}$.

\section{Rational-domain approximation with controlled denominators}
Here is the neighborhood point needed below. The ordinary topology of
an affinoid spectrum is the initial topology of its absolute-value
functions. Around a point in the interior of a polyannulus, choose
finitely many strict lower and upper bounds defining a neighborhood
contained in the prescribed open set. Move the endpoints slightly
inwards, keeping the point inside. Bounds of the form $|s|^{m/n}$ are
dense among positive real bounds, so every chosen constraint can be
written using an integral power of its function and a scalar in $K$.
Weierstrass upper constraints have denominator one. For the finitely
many positive lower constraints $|h_i|\ge a_i$, use common denominator
$h_0=\prod_i h_i$, numerators $h_0/h_i$, and add numerator one with
the redundant bound $1\le(\prod_i a_i)^{-1}|h_0|$. The added numerator
ensures the unit-ideal condition and excludes every zero of $h_0$;
on its complement the other inequalities are exactly the lower
constraints. Upper constraints multiply their numerators by $h_0$.
This gives a rational domain contained in the chosen neighborhood and
having the point in its ordinary interior. Integral powers realize
the scalar bounds in the ordinary strict presentation when needed.
Thus shrinking to such a compact rational neighborhood is justified
in the ordinary topology, not merely by a rigid admissible cover.

\begin{lemma}\label{rational}
Let $V$ be a rational domain in a compact torus polyannulus $P$, and
$F\in\mathcal O(V)$. There are rational functions
$F_n=A_n/Q_n$, with Laurent-polynomial numerators and denominators,
such that $Q_n$ is a unit on $V$ and $F_n\to F$ uniformly on $V$.
\end{lemma}
\begin{proof}
Write the defining presentation as
\[
 \mathcal O(V)=\mathcal O(P)\langle r_i^{-1}T_i\rangle/
                    (g_0T_i-g_i),
\]
where $g_0,g_1,\ldots,g_m$ generate the unit ideal and the domain is
defined by $|g_i|\le r_i|g_0|$. These inequalities and the unit-ideal
condition imply that $g_0$ has no zero on $V$. It is a unit there and
has a positive lower absolute-value bound by compactness.
Finite polynomials in $T_i$ with coefficients in $\mathcal O(P)$
approximate $F$. Each such expression is a finite sum of products of
the $g_i/g_0$ and finitely many coefficients. Laurent polynomials are
dense in the polyannulus algebra by its completed-series definition.
Approximate these finitely many coefficients and all $g_i$ on $P$.
Choose the approximation $q$ to $g_0$ with
$\|q-g_0\|_V<\inf_V|g_0|$. Then $|q|=|g_0|$ on $V$, and $q$ is a
unit there. Substitution in the fixed finite expression is continuous
in the uniform norm, because inverse denominators have this fixed
positive lower bound. Clearing a finite power of $q$ gives a quotient
$A/q^m$ of Laurent polynomials. Choosing each finite truncation and
its substitutions with error tending to zero proves the assertion.
No uniform bound on the degrees or on the denominator powers is needed.
\end{proof}

\section{All local principal analytic divisor tests}
\begin{theorem}\label{localdiv}
Let $u,v\in\mathcal C$ satisfy $\MA(u)=\MA(v)$, and let $W$ be an
ordinary open subset of a torus chart. For every regular holomorphic
function $F$ on $W$,
\[
 \DD\bigl((u-v)\delta_{\operatorname{div}F}\bigr)=0
                      \quad\text{on }W.
\]
This includes local analytic hypersurfaces not defined by a Laurent
polynomial on the entire torus.
\end{theorem}
\begin{proof}
Shrink around any point to the interior of a compact rational domain
$V\subset W$ in a polyannulus. Lemma\ref{rational} gives
$F_n=A_n/Q_n\to F$ uniformly. After omitting initial terms the numerator
is nonzero. A nonzero Laurent polynomial is regular on the smooth
torus, and $Q_n$ is a unit on $V$. Poincare--Lelong therefore gives the
positive current identity on its interior
\[
 T_n=\DD[\log|F_n|]
       =\delta_{\operatorname{div}A_n}.
\]
The preceding section gives $\DD(dT_n)=0$ there, where $d=u-v$.

For any compact smooth $(2,2)$ test $\alpha$ in this interior, its
support avoids the lower-dimensional zero locus of the regular $F$.
Hence $m_\alpha=\inf_{\operatorname{supp}\alpha}|F|>0$. Uniform
approximation implies $|F_n-F|<m_\alpha$ on that support for all late
$n$. The ultrametric inequality gives $|F_n|=|F|$ there. The currents
$[\log|F_n|]$ therefore agree with $[\log|F|]$ on each such test for
all sufficiently late $n$, and converge weakly. Differentiation and
Poincare--Lelong yield
\[
 T_n\longrightarrow\DD[\log|F|]
                          =\delta_{\operatorname{div}F}.
\]
These are positive currents. CLD8.5.1 now gives
$dT_n\to d\delta_{\operatorname{div}F}$, because $d$ is bounded
continuous. Differentiating and using $\DD(dT_n)=0$ proves the result
on the interior of $V$. The local assertions glue on $W$.
\end{proof}
This passage uses approximation on a rational neighborhood and
invertibility of its denominators. It makes no assertion that an
arbitrary analytic curve map has a local closed hypersurface image,
or that the current identity automatically survives every scalar
extension. The weighting step is not licensed for an unbounded
generated difference without a separate integrability theorem.


\section{From a divisor-current identity to a curve Laplacian}
We work first in the torus cover $Y=(\mathbb G_m^2)_K^{\rm an}$.
Let $d$ denote the pullback of the bounded continuous difference $u-v$.
The quotient map is a local analytic isomorphism. The divisor identities
proved above therefore hold on the ordinary open torus atlas of $Y$.

Let $P\in K[T]$ be a nonconstant irreducible polynomial with nonzero
constant term. Characteristic zero makes $P$ separable. The smooth,
reduced principal curve
\[
 C_P=\{P(z_1)=0\}\subset Y
\]
is $\mathbb G_{m,L}^{\rm an}$ viewed as a $K$-analytic space, where
$L=K[T]/(P)$. Its closed immersion $j$ is proper. The divisor current
of $P(z_1)$ is exactly $j_*\delta_{C_P}$, with multiplicity one.
For every chartwise continuous dpsh function $d$, continuous
multiplication and pushforward give
\[
 d\,j_*\delta_{C_P}=j_*((j^*d)\delta_{C_P}),\qquad
 \DD(d\delta_{\operatorname{div}P(z_1)})
       =j_*\DD_{C_P}[j^*d].                 \tag{1}
\]
Here $d$ is locally a difference of continuous psh functions because
the smooth reference potential cancels; pullback preserves this
minimal sheaf. The whole continuous projection-formula argument is
CLD8.5.6, extended by linearity; commutation of pushforward and both
differentials is CLD7.7.8. The same identity follows directly by
testing the weighted integration currents, so it does not require
closedness of a continuously weighted divisor current.

On the curve, $\DD_{C_P}[j^*d]$ is locally a difference of positive
scalar currents and thus a signed Radon measure. A proper closed
embedding pushes such a measure by the ordinary measure pushforward.
That pushforward is injective: a Borel set of the closed subspace is
a Borel subset of the ambient space. CLD7.2.8 identifies signed
Radon measures with their scalar currents, using density of compact
smooth functions. Equation(1) and the preceding divisor identity give
\[
                    \DD_{C_P}[d|_{C_P}]=0.                 \tag{2}
\]
This local conclusion glues along the whole noncompact curve.
The same assertion holds for the horizontal curve $z_2=1$.

There is a normalization to record before using(2) over $L$.
The underlying analytic functions and moment maps on $C_P$ are the
same when it is viewed over $K$ or over $L$. CLD14.4.3 states and
proves that its $K$-calibration is $[L:K]$ times its $L$-calibration.
Consequently the corresponding integration currents and Laplacians
have this same positive factor. Vanishing in(2) therefore gives
vanishing for the $L$-analytic curve with its own normalization.
We do not identify these two calibrations without their factor.

\section{Finite Galois base change of a zero signed Laplacian}
\begin{lemma}\label{finitegalois}
Let $L$ be a complete nontrivially valued field of characteristic zero,
let $E/L$ be a finite Galois extension, and let
$\pi:\mathbb G_{m,E}^{\rm an}\longrightarrow\mathbb G_{m,L}^{\rm an}$
be the base-change map. If $h$ is continuous, locally dpsh, and
$\DD[h]=0$, then $\DD[h\circ\pi]=0$ with the canonical
$E$-analytic normalization.
\end{lemma}
\begin{proof}
Pullback preserves continuous psh functions: the assertion holds for
smooth psh generators, pullback preserves locally uniform convergence,
and the inverse-image sheaf minimality proves it at every stage.
Thus $h\circ\pi$ is locally dpsh and
$\mu=\DD[h\circ\pi]$ is a signed Radon measure. It is invariant under
$G=\operatorname{Gal}(E/L)$, since $h\circ\pi$ is invariant and
the isometric automorphisms preserve the canonical calibration.

With the canonical normalization over the varying base field,
CLD15.3.1 and16.1.1 give $\pi_*\delta_{\mathbb G_{m,E}}
=\delta_{\mathbb G_{m,L}}$. Bounded continuous multiplication
locally on a compact test support therefore gives
\[
 \pi_*[h\circ\pi]=[h],\qquad
 \pi_*\mu=\DD\pi_*[h\circ\pi]=\DD[h]=0.        \tag{3}
\]
For clarity, if the source is instead viewed as an $L$-analytic
finite cover, its calibration is multiplied by $[E:L]$ and both
equalities acquire that factor. Either convention proves the same
vanishing. We use the varying-base convention in(3).

Every fibre of $\pi$ is a single $G$-orbit. Indeed it is the spectrum
of the finite etale algebra $\mathscr H(x)\otimes_L E$, a product
of finite field extensions of the complete field $\mathscr H(x)$.
The Galois group permutes its factors transitively: its invariant
algebra is $\mathscr H(x)$, which has no nontrivial idempotent.
Each factor has a unique extending absolute value and contributes
one analytic point. The finite proper surjective map $\pi$ is closed,
hence a topological quotient.

For $\varphi\in C_c(\mathbb G_{m,E}^{\rm an})$ put
$\overline\varphi=|G|^{-1}\sum_{\sigma\in G}\varphi\circ\sigma$.
It is continuous, compactly supported, and constant on each fibre,
so it descends to $\psi\in C_c(\mathbb G_{m,L}^{\rm an})$.
Invariance of $\mu$ and(3), interpreted as signed measures through
CLD7.2.8, imply
\[
 \int\varphi\,d\mu=\int\overline\varphi\,d\mu
                      =\int\psi\,d(\pi_*\mu)=0.
\]
Hence $\mu=0$. This proves only the stated finite Galois passage;
no unrestricted pullback theorem for distributional currents is used.
\end{proof}

\section{Gauss-ray Liouville lemma}
\begin{lemma}\label{liouville}
Let $L$ be a complete nontrivially valued field of characteristic zero.
A bounded continuous locally dpsh function $h$ on
$\mathbb G_{m,L}^{\rm an}$ with $\DD[h]=0$ is constant.
\end{lemma}
\begin{proof}
Write $\eta_{a,r}$ for the Gauss point of the closed disk centred at
$a$ with radius $r>0$; its polynomial norm is
$\left|\sum b_i(z-a)^i\right|=\max_i|b_i|r^i$.
Put $H_0(t)=h(\eta_{0,e^{-t}})$ for $t\in\mathbb R$.
For $\phi\in C_c^\infty(\mathbb R)$ the function
$\phi(-\log|z|)$ is smooth with compact support in a closed annulus
inside $\mathbb G_m$. The canonical one-dimensional torus calibration
and integration give
\[
 0=\langle\DD[h],\phi(-\log|z|)\rangle
       =\int_{\mathbb R} H_0(t)\phi''(t)\,dt.       \tag{4}
\]
This is the ordinary calibrated torus moment $z$, of degree one:
$\DD\phi(t)=\phi''(t)d't\wedge d''t$ and its scalar measure is
the pushforward of $\phi''(t)dt$ by the Gauss section. This follows
directly by integrating the one-dimensional calibrated cell; there
is no multiplicity from a larger ground field in the $L$-normalization.
The section and its norm formula are CLD12.1.2,12.1.5.
Equation(4) makes $H_0$ affine as a distribution on $\mathbb R$.
Continuity identifies it with that affine function; boundedness
forces its slope to vanish. Denote its constant value by $c$.

Take $a\in L^*$. On the punctured disk $0<|z-a|<|a|$ the coordinate
$z-a$ is a unit and the disk is contained in $\mathbb G_m$.
The same degree-one annulus calculation, with compact test support,
makes
\[
 H_a(t)=h(\eta_{a,e^{-t}}),\qquad t>-\log|a|,
\]
affine on this half-line. Boundedness as $t\to+\infty$ forces slope
zero. As $r\downarrow0$, $\eta_{a,r}$ converges to the rational
point $a$; this follows from the displayed polynomial norm and the
initial seminorm topology. As $r\uparrow|a|$, it converges to
$\eta_{a,|a|}=\eta_{0,|a|}$. The equality follows by expanding
polynomials at either centre and applying the ultrametric inequality
in both directions. Continuity now gives $h(a)=c$.

A closed rigid point $x$ of $\mathbb G_{m,L}^{\rm an}$ has finite
residue field over $L$. Choose a finite Galois normal closure $E/L$
of that residue field; characteristic zero supplies separability.
Lemma\ref{finitegalois} applies to $h_E=h\circ\pi$. Its Gauss
section at zero projects to the original Gauss section by the norm
formula (also CLD12.1.8(a)). Thus the constant in its Gauss-ray
calculation is the same $c$. A lift of $x$ is an $E$-rational point,
and the previous paragraph proves $h(x)=c$.

We justify the density used to finish the proof. Every nonempty
ordinary open subset of the strict analytic torus contains a
nonempty strict rational affinoid domain: move finite absolute-value
bounds inward while keeping a chosen point inside and replace their
endpoints by elements of $|L^*|^{\mathbb Q}$; that group is dense
since the valuation is nontrivial. Integral powers realize these
constraints with scalars in $L$, using the common-denominator
unit-ideal construction given above. A compact polyannulus containing
the construction can likewise be given a strict presentation by
integral powers. The resulting nonzero strict affinoid algebra has
a maximal ideal. The Tate Nullstellensatz gives a finite residue
extension of $L$, with its unique absolute value, and therefore a
closed rigid analytic point in that domain. Here the Nullstellensatz
is the named ordinary affinoid foundation; Van der Put1.6 records
the exact finite-residue statement, and Papikian3.2(d) corroborates it.
Their short immediate passages do not contain a whole proof of
Weierstrass preparation. No fresh proof of that deeper result is claimed.
Closed rigid points are therefore dense. Since $h$ is continuous and
equals $c$ on them, it is identically $c$.
\end{proof}

\section{Coordinate fibres prove bounded uniqueness}
\begin{theorem}\label{boundedunique}
On the compact analytic quotient specified in the first section, let
$u,v\in\mathcal C$. If their actual recursively defined measures
$\MA(u)$ and $\MA(v)$ agree, then $u-v$ is constant pointwise on $X$.
\end{theorem}
\begin{proof}
The difference pulled back to $Y$ is bounded and continuous.
Equations(1)--(2) and the calibration comparison show that its
restriction to every vertical $C_P\cong\mathbb G_{m,L}^{\rm an}$
is locally continuous dpsh with zero Laplacian over $L$.
Lemma\ref{liouville} makes each such restriction constant.
The same lemma on the horizontal $K$-curve $z_2=1$ makes its
restriction a single constant $c$ there.

Let $y$ be any closed rigid point of $Y$. Its coordinates have values
$(a,b)$ in a finite extension $E/K$. Let $P$ be the minimal polynomial
of $a$ over $K$ and $L=K(a)$. Then $y$ lies on $C_P$.
The point $(a,1)$ of that curve lies also on $z_2=1$ and has value
$c$. Constancy on $C_P$ therefore gives $d(y)=c$, irrespective of
whether $b$ lies in $L$ or in a larger finite extension. This uses
the entire $L$-analytic curve, not just its $L$-rational points.

Closed rigid points are dense in $Y$ by the same strict rational
domain and Nullstellensatz argument in dimension two.
Continuity gives $d=c$ on all of $Y$. The surjective quotient map
then gives $u-v=c$ on all of $X$.
\end{proof}
The proof of Theorem\ref{boundedunique} needs only the Laurent
divisor tests, although the full local holomorphic divisor extension
has also been retained above. General curve recognition, normalization
of arbitrary curve images, and infinite scalar extension are not
premises of the coordinate-fibre proof.

\section{Exact reading scope and original residual}
The whole mathematical body of the earlier working paper398 through
its local analytic divisor theorem has been included above. The new
steps are the injective scalar-measure passage on coordinate curves,
the finite Galois averaging argument with the explicit normalization,
the Gauss-ray calculation, and the pointwise coordinate-fibre connection.
All of these working claims are UNREVIEWED.

For CLD the retained whole v2 body has SHA256
\texttt{64846cb00926267c6238d1afac61b5eaed0a0e3a6349725738f9b1bf163de085}.
The entire immediate applied8.5.6 projection proof,7.2.8 scalar measure
proof,7.7.8 differential commutation passage,14.4.3 calibration
comparison,15.3.1 base-change integration proof,16.1.1 current proof,
and12.1.2/5/8(a) Gauss norm passages were freshly read in bounded
literal-LF chunks. The whole immediate Poincare--Lelong statement
and proof16.7.3 were already freshly read for398. Whole primary source
bodies, exact literal units and actual pixels remain local; the public
record reports their bindings and qualifications. Whole8.6.7 was read
but its arbitrary continuous multiplier closedness is not a premise.
The Van der Put actual original PDF page3 (printed J7-02), including
the finite-residue statement1.6 and its immediate reference to1.2,
was freshly viewed; the entire deeper normalization proof1.2 is not
claimed newly audited. Its pages5--6, containing the whole2.2 rational
presentation argument, were freshly viewed for398. Papikian3.2 is
corroborating statement evidence, with no complete immediate proof
printed there. Bounded proofreading is not an audit of every recursive
source dependency. Ordinary affinoid foundations remain explicit.

Bounded continuous uniqueness is the scope of
Theorem\ref{boundedunique}. Original P16.7 still requires much more:
intrinsic recognition of unbounded generated finite-energy functions,
integration and pointwise restrictions outside bounded strata,
actual compactness of finite-energy sublevels, obstacle envelopes and
orthogonality, existence for every eligible arbitrary-support measure,
and uniqueness for the full unbounded equation domain. None of these
is established by a Liouville lemma whose proof uses boundedness.
No original catalogue item is closed and no accepted-paper count is
increased. New papers398--399 were not supplied as hot feedback to
the independently running native278--279 campaigns; their frozen
inputs and criteria remain unchanged.

\begin{thebibliography}{9}
\bibitem{CLD} A. Chambert-Loir and A. Ducros,
\emph{Formes differentielles reelles et courants sur les espaces de
Berkovich}, arXiv:1204.6277v2, exact retained source snapshot.
\bibitem{GR} W. Gubler and J. Rabinoff,
\emph{Classically psh and pluriharmonic functions on Berkovich spaces},
arXiv:2506.13548v2, 17 September 2025.
\bibitem{VdP} M. van der Put, \emph{Rigid analytic spaces},
Groupe de travail d'analyse ultrametrique3(2), exp.J7 (1975--1976),
Sections1.1 and2.1--2.2, exact retained Numdam scan.
\bibitem{Pap} M. Papikian, \emph{Rigid analytic geometry and abelian
varieties}, notes dated13 June2004, Sections3.1--3.2, Definition3.4
and Theorem3.5.
\end{thebibliography}
\end{document}
